\scanpage{59}
\textcircled{A} $\mathcal K$ catégorie avec $\varprojlim$ finies

$\mathcal M\xrightarrow{\varphi}\mathcal K$
$\left\{\begin{array}{l}
\text{foncteur exact à g. (avec }\varprojlim\text{ finies existant dans }\mathcal M)\\
\text{conservatif (donc fidèle)}\\
I_0\text{ objet de }\mathcal M\text{ tel que }\forall X\in\operatorname{Ob}\mathcal M,\ \exists n\in\mathbb N
\text{ et mono }i:X\longrightarrow I_0^n\\
\deleted{\varphi\text{ transportable}}
\end{array}\right.$

À équivalence près (pour $(\mathcal M,\varphi)$), pour $\mathcal K$ donné, ces données équivalent à la donnée
\begin{enumerate}
\item[a)] d'un $I\in\operatorname{Ob}\mathcal K$
\item[b)] $\forall n\in\mathbb N$, d'un ens $M_n$ de sous-objets de $I^n$
\end{enumerate}
avec les conditions

$1^\circ)$ $M_n$ stable par intersections finies

$2^\circ)$ $M_n$ stable par image inverse d'appl. simpl. $I^n\to I^{n'}$

$3^\circ)$ Pour $\deleted{\sigma:I^n\to I^{n'}\ \unsure{\text{comme ci-dessus}}}$ \inserted{appl.}, et $K\in M_n$ tel que $\sigma|K$ soit mono, $\sigma(K)\in M_{n'}$.

\uline{NB} Si on \uline{suppose} $\varphi$ transportable, on a correspondance avec $\mathcal M$ défini à $\mathcal K$-\uline{isom} près.

\textcircled{B} Supposons que tt morphisme $X\xrightarrow{f}Y$ de $\mathcal K$ se factorise en $X\xrightarrow{g}Y'\lhook\joinrel\xrightarrow{i}Y$, $g$ épi effectif, $i$ monomorphisme (donc on peut écrire $Y'=\operatorname{Im}g$, c'est le plus petit sous-objet de $Y$ par lequel se factorise $g$).
\inserted{\diagram{\begin{tikzcd}[ampersand replacement=\&]
X \arrow[r,dash,"{g}"] \arrow[drr] \& Y' \arrow[r,hook,"{i}"] \arrow[dr] \& Y \arrow[d,"{\unsure{i'|\sim}}"]\\
\& \& Y''
\end{tikzcd}}}
\inserted{Supposons de plus les épi effectifs le soient universellement.}
Considérons la condition plus forte que $3^\circ$

$(3^\circ\text{bis})$ Comme $3^\circ$, mais sans supposer que $\sigma|K$ soit mono.

Cette condition est équivalente à : tout sous-objet d'un objet de $\mathcal K$ à ``structure modérée'', image d'un morphisme ``modéré'' $f$, soit modéré.

\margin{$\Longleftrightarrow$ Tout $f:X\to Y$ dans $\mathcal M$ se factorise en $X\xrightarrow{g}Y'\lhook\joinrel\xrightarrow{i}Y$, $i$ mono, $g$ \deleted{\unsure{universellement}} épi effectif. \deleted{De} plus $\varphi$ transforme épi eff. en \unsure{tels}.
(En effet, on aura : $\varphi(g)$ épi eff. ssi $g$ épim. eff., \unsure{d'où} les épi eff. de $\mathcal M$ sont univers.)
$\Longrightarrow$ Si $g:X\to Y$ épi eff. dans $\mathcal M$, et $Z\in\operatorname{Ob}\mathcal M$, $h:Y\to Z$ dans $\mathcal K$, tel que $h\circ g$ modéré, alors $h$ modéré.}

\textcircled{C} Si dans $\mathcal K$ les sommes finies existent \inserted{(et sont disjointes, et universelles)}, dans \inserted{l'ens} des sous-objets d'un objet de $\mathcal K$, il y a des sup finis (``réunion''). Considérons

$4^\circ)$ a) $M_n$ stable par \deleted{sup} \inserted{sup} finis $(\forall n\geq0)$ \deleted{$\unsure{\varnothing}$}

b) Il existe un \unsure{morphisme} modéré de \deleted{\unsure{deux éléments égaux}} $e\amalg e$ dans $I$ [défini par deux \inserted{$e_1,e_2$} $\in M_1$ tels que $e_1\overset{\sim}{\to}e$ et $e_2\overset{\sim}{\to}e$] qui soit un monomorphisme

\margin{60}
